\section{Methodology}
\label{sec:method}

This section presents the Deep Conformal Alignment (DCA) framework for learning topology-preserving, angle-preserving transformations that normalize morphological variability in histopathology images. We begin with an overview of the architecture and key innovations, then formalize the problem mathematically, detail the diffeomorphic transformation network, introduce the conformal regularization derived from differential geometry, and finally present the end-to-end learning objective with its optimization procedure.

\subsection{Overview and Key Innovations}
\label{subsec:overview}

Figure~\ref{fig:architecture} illustrates the DCA architecture. Unlike standard CNNs that directly process raw images, DCA introduces a learnable geometric normalization module that warps inputs into canonical representations before classification. The framework makes three key technical contributions that distinguish it from prior work.

First, while standard Spatial Transformer Networks (STN)~\cite{jaderberg2015spatial} learn unconstrained transformations that may introduce topological artifacts such as folding or tearing, we parameterize deformations via velocity fields integrated through scaling-and-squaring~\cite{arsigny2006logdiff}, guaranteeing diffeomorphism. This is critical for biomedical imaging where tissue connectivity must be preserved. Second, we introduce a differentiable loss derived from the Cauchy-Riemann equations in complex analysis, which constrains learned transformations to be approximately conformal, thereby preserving angles. This maintains local tissue structure integrity while normalizing global morphology. Our formulation connects to the Beltrami coefficient, providing theoretical grounding in quasiconformal mapping theory. Third, unlike prior geometric normalization methods~\cite{huang2024surface} that require expensive per-image optimization with $O(n^3)$ complexity for $n$ pixels, DCA amortizes this cost through learning, enabling real-time inference with $O(n)$ forward pass while adapting transformations to maximize task-specific performance.

Table~\ref{tab:comparison_prior} contrasts DCA with related approaches. Standard STN learns unconstrained affine or thin-plate spline transformations without geometric constraints. Fixed conformal mapping enforces conformality but requires per-image optimization and cannot adapt to semantic content. DCA uniquely combines learnable transformations with explicit geometric constraints, achieving both efficiency and mathematical rigor.

\begin{table}[t]
\centering
\caption{Comparison of Geometric Transformation Approaches}
\label{tab:comparison_prior}
\begin{tabular}{lccc}
\hline
\textbf{Method} & \textbf{Learnable} & \textbf{Conformal} & \textbf{Diffeomorphic} \\
\hline
Standard STN~\cite{jaderberg2015spatial} & \checkmark & $\times$ & $\times$ \\
Fixed CEM~\cite{huang2024surface} & $\times$ & \checkmark & \checkmark \\
\textbf{DCA (Ours)} & \checkmark & \checkmark & \checkmark \\
\hline
\end{tabular}
\end{table}

\subsection{Problem Formulation}
\label{subsec:formulation}

Let $\Omega \subset \mathbb{R}^2$ denote the image domain, typically $[0,H] \times [0,W]$ for an $H \times W$ image, and let $\mathcal{I} = L^2(\Omega, \mathbb{R}^3)$ denote the space of RGB images. Given a labeled dataset $\mathcal{D} = \{(I_i, y_i)\}_{i=1}^N$ where $I_i \in \mathcal{I}$ and $y_i \in \{1, \ldots, C\}$ denotes the tissue class, we seek to learn two functions jointly. The first is a transformation network $\mathcal{T}_\psi: \mathcal{I} \rightarrow \mathfrak{X}(\Omega)$ that maps each image to a velocity field $v_i = \mathcal{T}_\psi(I_i)$, where $\mathfrak{X}(\Omega)$ denotes the Lie algebra of smooth vector fields on $\Omega$. From this velocity field, we obtain a diffeomorphism $\phi_i = \exp(v_i) \in \text{Diff}^+(\Omega)$, where $\text{Diff}^+(\Omega)$ denotes the group of orientation-preserving diffeomorphisms. The second function is a classifier $f_\theta: \mathcal{I} \rightarrow \Delta^{C-1}$ that predicts class probabilities, where $\Delta^{C-1}$ is the $(C-1)$-dimensional probability simplex. Table~\ref{tab:notation} summarizes the key mathematical notation used throughout this section.

\begin{table}[t]
\centering
\caption{Mathematical Notation}
\label{tab:notation}
\begin{tabular}{ll}
\hline
\textbf{Symbol} & \textbf{Description} \\
\hline
$\Omega \subset \mathbb{R}^2$ & Image domain \\
$I: \Omega \rightarrow \mathbb{R}^3$ & RGB image \\
$v: \Omega \rightarrow \mathbb{R}^2$ & Velocity field \\
$\phi: \Omega \rightarrow \Omega$ & Diffeomorphism (deformation) \\
$J_\phi \in \mathbb{R}^{2 \times 2}$ & Jacobian matrix of $\phi$ \\
$\text{Diff}^+(\Omega)$ & Group of orientation-preserving diffeomorphisms \\
$\mathfrak{X}(\Omega)$ & Lie algebra of vector fields \\
$\mathcal{E}_{conf}(\phi)$ & Conformal energy \\
\hline
\end{tabular}
\end{table}

The transformed image is obtained as $\tilde{I}_i = I_i \circ \phi_i$ by warping $I_i$ through $\phi_i$ via differentiable bilinear interpolation. We formulate the learning problem as a constrained optimization:
\begin{equation}
\min_{\psi, \theta} \quad \mathbb{E}_{(I,y) \sim \mathcal{D}} \left[ \mathcal{L}_{cls}(f_\theta(\tilde{I}), y) + \lambda_{conf} \mathcal{E}_{conf}(\phi) + \lambda_{smooth} \mathcal{R}_{smooth}(\phi) \right]
\label{eq:objective}
\end{equation}
where $\mathcal{L}_{cls}$ is the cross-entropy classification loss, $\mathcal{E}_{conf}(\phi)$ is the conformal energy measuring deviation from angle preservation, $\mathcal{R}_{smooth}(\phi)$ is a smoothness regularizer preventing unrealistic deformations, and $\lambda_{conf}, \lambda_{smooth} > 0$ are regularization weights. This objective balances three competing goals: maximizing classification accuracy, enforcing geometric constraints, and maintaining smooth transformations.

We make three key assumptions. First, the velocity field $v$ is sufficiently smooth ($v \in C^2(\Omega)$) to ensure $\phi = \exp(v)$ is a diffeomorphism. Second, an optimal canonical representation exists and is approximately conformal for each tissue class. Third, the classifier $f_\theta$ has sufficient capacity to learn discriminative features from normalized images. These assumptions are reasonable for histopathology images where tissue structures exhibit local regularity despite global variability.

\subsection{Diffeomorphic Spatial Transformer Network}
\label{subsec:diffeomorphic}

Directly parameterizing the infinite-dimensional Lie group $\text{Diff}^+(\Omega)$ is challenging. We adopt the Lie algebra approach~\cite{arsigny2006logdiff,voxelmorph} where instead of predicting $\phi$ directly, we predict a stationary velocity field $v: \Omega \rightarrow \mathbb{R}^2$ and obtain $\phi$ via the exponential map. This approach has strong theoretical foundations in differential geometry.

\begin{definition}[Exponential Map]
\label{def:exp_map}
The exponential map $\exp: \mathfrak{X}(\Omega) \rightarrow \text{Diff}^+(\Omega)$ is defined as:
\begin{equation}
\phi = \exp(v) = \lim_{T \rightarrow \infty} \left( \text{Id} + \frac{v}{T} \right)^T
\end{equation}
where $\text{Id}$ denotes the identity transformation.
\end{definition}

Let $g_\psi: \mathcal{I} \rightarrow L^2(\Omega, \mathbb{R}^2)$ be a convolutional neural network that we call the localization network. Given input $I$, we compute the velocity field as $v = g_\psi(I)$. We implement $g_\psi$ as a U-Net~\cite{ronneberger2015unet} with four downsampling blocks in the encoder, each containing two $3 \times 3$ convolutions followed by ReLU activation and $2 \times 2$ max pooling. The channel dimensions progress as $[3, 64, 128, 256, 512]$ through the encoder. The decoder mirrors this structure with four upsampling blocks using skip connections, where each block applies $2 \times 2$ transposed convolution followed by two $3 \times 3$ convolutions and ReLU. The output layer is a $1 \times 1$ convolution producing two channels corresponding to the velocity field components $v_x$ and $v_y$. Critically, we initialize the final layer weights and biases to zero, ensuring $v \approx 0$ at initialization, which corresponds to the identity transformation. This initialization strategy is essential for stable training.

To compute $\phi = \exp(v)$, we use the scaling-and-squaring algorithm~\cite{arsigny2006logdiff}, which approximates the exponential map via repeated self-composition. The algorithm first scales the velocity field by $v_0 = v / 2^T$ where $T$ is the number of squaring steps, then iteratively computes $\phi_t = \phi_{t-1} + \phi_{t-1} \circ \phi_{t-1}$ for $t = 1, \ldots, T$. The composition $\phi \circ \phi$ is implemented via differentiable bilinear interpolation: given $\phi(x) = (u(x), v(x))$, we evaluate $\phi(\phi(x)) = \phi(x + \phi(x))$ by sampling $\phi$ at location $x + \phi(x)$. Algorithm~\ref{alg:scaling_squaring} provides the complete procedure.

\begin{algorithm}[t]
\caption{Scaling-and-Squaring for Diffeomorphic Integration}
\label{alg:scaling_squaring}
\begin{algorithmic}[1]
\REQUIRE Velocity field $v \in \mathbb{R}^{H \times W \times 2}$, number of steps $T \in \mathbb{N}$
\ENSURE Deformation field $\phi \approx \exp(v)$
\STATE $\phi_0 \leftarrow v / 2^T$
\FOR{$t = 1$ to $T$}
    \STATE $\phi_t(x) \leftarrow \phi_{t-1}(x) + \phi_{t-1}(x + \phi_{t-1}(x))$ via bilinear interpolation
\ENDFOR
\RETURN $\phi_T$
\end{algorithmic}
\end{algorithm}

The theoretical guarantee for this approach is formalized in the following proposition.

\begin{proposition}[Diffeomorphism Guarantee]
\label{prop:diffeomorphism}
If $\|v\|_{L^\infty} < \epsilon$ for sufficiently small $\epsilon > 0$, then $\phi = \exp(v)$ obtained via Algorithm~\ref{alg:scaling_squaring} with $T \geq 7$ is a diffeomorphism with probability $> 1 - \delta$ for arbitrarily small $\delta > 0$.
\end{proposition}

\begin{proof}[Proof Sketch]
The result follows from Arsigny et al.~\cite{arsigny2006logdiff}: for small $v$, the exponential map $\exp: \mathfrak{X}(\Omega) \rightarrow \text{Diff}^+(\Omega)$ is a local diffeomorphism from the Lie algebra of vector fields. The scaling-and-squaring approximation converges exponentially fast with rate $\|\phi_T - \exp(v)\| = O(2^{-T})$. For $T=7$, the approximation error is less than $10^{-3}$ in practice. The condition $\|v\|_{L^\infty} < \epsilon$ ensures the Jacobian determinant $\det(J_\phi) > 0$ everywhere, guaranteeing orientation preservation.
\end{proof}

The computational complexity of Algorithm~\ref{alg:scaling_squaring} is $O(THW)$ where $H \times W$ is the image resolution, as it requires $T$ self-compositions, each involving bilinear interpolation over all pixels. This linear complexity in image size makes the approach practical for high-resolution histopathology images.

\subsection{Conformal Regularization via Differential Geometry}
\label{subsec:conformal}

A smooth map $\phi: \Omega \rightarrow \Omega$ is conformal if it preserves angles locally. In two dimensions, conformality is elegantly characterized by the Cauchy-Riemann equations from complex analysis. This geometric property is crucial for histopathology because it maintains local tissue structure (e.g., cell arrangements, gland orientations) while normalizing global morphology.

\begin{definition}[Conformal Map]
\label{def:conformal}
Let $\phi(x,y) = (u(x,y), v(x,y))$ be a smooth map with Jacobian matrix:
\begin{equation}
J_\phi = \begin{bmatrix} \frac{\partial u}{\partial x} & \frac{\partial u}{\partial y} \\ \frac{\partial v}{\partial x} & \frac{\partial v}{\partial y} \end{bmatrix}
\end{equation}
The map $\phi$ is conformal if and only if $J_\phi$ is a scaled rotation:
\begin{equation}
J_\phi = \lambda(x,y) \begin{bmatrix} \cos\theta(x,y) & -\sin\theta(x,y) \\ \sin\theta(x,y) & \cos\theta(x,y) \end{bmatrix}
\label{eq:conformal_jacobian}
\end{equation}
for some scaling factor $\lambda > 0$ and rotation angle $\theta$. This is equivalent to the Cauchy-Riemann equations:
\begin{equation}
\frac{\partial u}{\partial x} = \frac{\partial v}{\partial y}, \quad \frac{\partial u}{\partial y} = -\frac{\partial v}{\partial x}
\label{eq:cauchy_riemann}
\end{equation}
\end{definition}

Since our transformation is $F(x) = x + \phi(x)$ in displacement formulation, we measure conformality of $F$ whose Jacobian is $J_F = I + J_\phi$. We define the conformal energy as the $L^2$ deviation from the Cauchy-Riemann equations:
\begin{equation}
\mathcal{E}_{conf}(\phi) = \frac{1}{|\Omega|} \int_\Omega \left[ \left( \frac{\partial u}{\partial x} - \frac{\partial v}{\partial y} \right)^2 + \left( \frac{\partial u}{\partial y} + \frac{\partial v}{\partial x} \right)^2 \right] dx
\label{eq:conformal_energy}
\end{equation}
This can be equivalently written using the Frobenius norm as $\mathcal{E}_{conf}(\phi) = \frac{1}{|\Omega|} \int_\Omega \| J_\phi - \frac{\text{tr}(J_\phi)}{2} I \|_F^2 dx$ where $\text{tr}(J_\phi) = \frac{\partial u}{\partial x} + \frac{\partial v}{\partial y}$ is the trace. In practice, we discretize using central finite differences: $\frac{\partial u}{\partial x}(i,j) \approx \frac{u(i+1,j) - u(i-1,j)}{2\Delta x}$ and similarly for other partial derivatives. The discrete conformal energy becomes:
\begin{equation}
\mathcal{E}_{conf}(\phi) = \frac{1}{HW} \sum_{i,j} \left[ \left( \nabla_x u_{ij} - \nabla_y v_{ij} \right)^2 + \left( \nabla_y u_{ij} + \nabla_x v_{ij} \right)^2 \right]
\label{eq:discrete_conformal}
\end{equation}

Our formulation connects to the Beltrami coefficient $\mu: \Omega \rightarrow \mathbb{C}$, a complex-valued measure of quasiconformality defined as $\mu = \frac{\partial_{\bar{z}} \phi}{\partial_z \phi}$ where $z = x + iy$ and $\partial_z = \frac{1}{2}(\partial_x - i\partial_y)$, $\partial_{\bar{z}} = \frac{1}{2}(\partial_x + i\partial_y)$. The following proposition establishes this connection.

\begin{proposition}[Conformality and Beltrami Coefficient]
\label{prop:beltrami}
A map $\phi$ is conformal if and only if $\mu = 0$ almost everywhere. Furthermore, the conformal energy satisfies:
\begin{equation}
\mathcal{E}_{conf}(\phi) = \frac{4}{|\Omega|} \int_\Omega |\mu|^2 \left| \frac{\partial \phi}{\partial z} \right|^2 dx
\end{equation}
\end{proposition}

\begin{proof}
From the definition of $\mu$ and the Cauchy-Riemann equations, we have $\mu = 0 \Leftrightarrow \partial_{\bar{z}} \phi = 0 \Leftrightarrow$ Eq.~\eqref{eq:cauchy_riemann} holds. The energy equivalence follows from direct computation using $\partial_z \phi = \frac{1}{2}[(\partial_x u - \partial_y v) + i(\partial_x v + \partial_y u)]$ and expanding the squared terms.
\end{proof}

This connection provides theoretical grounding in quasiconformal mapping theory: minimizing $\mathcal{E}_{conf}$ is equivalent to minimizing the $L^2$ norm of the Beltrami coefficient, which measures deviation from conformality. This allows us to interpret our learned transformations through the lens of established mathematical theory.

\subsection{End-to-End Learning Objective}
\label{subsec:objective}

The full training objective combines three terms that balance classification performance, geometric constraints, and transformation smoothness:
\begin{equation}
\mathcal{L}(\psi, \theta; I, y) = \mathcal{L}_{cls}(f_\theta(\tilde{I}), y) + \lambda_{conf} \mathcal{E}_{conf}(\phi) + \lambda_{smooth} \mathcal{R}_{smooth}(\phi)
\label{eq:full_loss}
\end{equation}
where $\mathcal{L}_{cls} = -\sum_{c=1}^C y_c \log f_\theta(\tilde{I})_c$ is the cross-entropy loss, $\mathcal{E}_{conf}(\phi)$ is given by Eq.~\eqref{eq:discrete_conformal}, and $\mathcal{R}_{smooth}(\phi) = \frac{1}{|\Omega|} \sum_{x \in \Omega} \|\nabla \phi(x)\|_F^2$ is the total variation smoothness regularizer. We set $\lambda_{conf} = 1.0$ and $\lambda_{smooth} = 0.1$ based on grid search over $\{0.1, 0.5, 1.0, 5.0\} \times \{0.01, 0.1, 1.0\}$ on the validation set. The conformal weight $\lambda_{conf}$ controls the trade-off between task performance and geometric constraint satisfaction.

We optimize Eq.~\eqref{eq:full_loss} using Adam~\cite{kingma2014adam} with initial learning rate $\alpha_0 = 10^{-3}$ and exponential decay factor $0.95$ applied every 10 epochs. Algorithm~\ref{alg:training} presents the complete training procedure. For each mini-batch, we first predict velocity fields using the localization network, then apply scaling-and-squaring to obtain diffeomorphisms, warp the images via bilinear interpolation, classify the transformed images, and finally update parameters via backpropagation. All components are differentiable, enabling end-to-end gradient flow from the classification loss through the geometric transformations back to the velocity field predictor.

\begin{algorithm}[t]
\caption{DCA Training Procedure}
\label{alg:training}
\begin{algorithmic}[1]
\REQUIRE Dataset $\mathcal{D} = \{(I_i, y_i)\}_{i=1}^N$, batch size $B$, epochs $E$, learning rate $\alpha_0$
\ENSURE Trained parameters $\psi^*, \theta^*$
\STATE Initialize $\psi, \theta$ randomly (final layer of $g_\psi$ initialized to zero)
\STATE Set learning rate $\alpha \leftarrow \alpha_0 = 10^{-3}$
\FOR{epoch $e = 1$ to $E$}
    \FOR{each mini-batch $\{(I_b, y_b)\}_{b=1}^B \sim \mathcal{D}$}
        \STATE $v_b \leftarrow g_\psi(I_b)$
        \STATE $\phi_b \leftarrow \text{ScalingSquaring}(v_b, T=7)$
        \STATE $\tilde{I}_b \leftarrow I_b \circ \phi_b$ via bilinear interpolation
        \STATE $\hat{y}_b \leftarrow f_\theta(\tilde{I}_b)$
        \STATE Compute $\mathcal{L}(\psi, \theta; I_b, y_b)$ via Eq.~\eqref{eq:full_loss}
        \STATE $(\psi, \theta) \leftarrow \text{Adam}(\nabla_{\psi,\theta} \mathcal{L}, \alpha)$
    \ENDFOR
    \IF{$e \mod 10 = 0$}
        \STATE $\alpha \leftarrow \alpha \times 0.95$
    \ENDIF
\ENDFOR
\RETURN $\psi^*, \theta^*$
\end{algorithmic}
\end{algorithm}

The per-image computational cost is $O(HW(D + D' + T))$ where $H \times W$ denotes image resolution, $D$ is the depth of the localization network $g_\psi$, $D'$ is the depth of the classifier $f_\theta$, and $T=7$ is the number of scaling-and-squaring steps. During training, we compute gradients through all components with this complexity per image. During inference, we perform a single forward pass with the same complexity. Critically, unlike fixed conformal mapping methods~\cite{huang2024surface} that require $O(n^3)$ per-image optimization where $n = HW$, DCA amortizes this cost during training, enabling real-time deployment. We implement the framework in PyTorch 1.12 with CUDA 11.3 on an NVIDIA A100 GPU with 40GB memory. Using batch size 16 for $224 \times 224$ images, training takes approximately 2 hours for 10 epochs on 5,026 training images, while inference requires only 15ms per image including both transformation and classification. The classifier uses ResNet-18~\cite{he2016resnet} pretrained on ImageNet and fine-tuned end-to-end with the transformation network.
